\section{Macros and Customisation}\label{sec:macros}

Many converters treat a \LaTeX{} document as a tree of known commands and give up on
anything else. \texdocx{} instead rebuilds \TeX's front end: a tokenizer with category
codes and a macro expander with grouping, conditionals, registers and counters. Your own
definitions are therefore \emph{executed}, the way \LaTeX{} would execute them, and
the Word file receives their result rather than the raw macro name. Every example
macro below is defined in this section's own source file and then used live; the code
blocks show the definitions verbatim.

\subsection{Commands with optional arguments}

\cmd{newcommand} accepts an optional first argument with a default value, just like in
\LaTeX. The following macro typesets a keyword in bold and colours it, using the accent
colour from the preamble unless another colour is given:

\begin{verbatim}
\newcommand{\keyword}[2][accent]{\textcolor{#1}{\textbf{#2}}}
\end{verbatim}
\newcommand{\keyword}[2][accent]{\textcolor{#1}{\textbf{#2}}}

Writing \verb|\keyword{expansion}| gives \keyword{expansion}, while
\verb|\keyword[accent!50]{tint}| gives \keyword[accent!50]{tint}. The second form uses an
\texttt{xcolor} mixing expression: \texttt{accent!50} is half accent and half white, and
\texttt{accent!60!black} would mix the accent with black instead. \texdocx{} evaluates
these expressions itself and writes the resulting RGB value into the run, so
\textcolor{accent}{plain accent text} and \textcolor{accent!50}{the lighter mix} both
survive as ordinary Word font colours that you can change later.

\subsection{Environments}

\cmd{newenvironment} works with arguments and an optional default, and the end code runs
before the environment name is checked, as in \LaTeX:

\begin{verbatim}
\newenvironment{aside}[1][Aside]
  {\par\medskip\noindent\textcolor{accent}{\textbf{#1.}}\ \itshape}
  {\par\medskip}
\end{verbatim}
\newenvironment{aside}[1][Aside]
  {\par\medskip\noindent\textcolor{accent}{\textbf{#1.}}\ \itshape}
  {\par\medskip}

\begin{aside}[Why it matters]
Declarations such as \cmd{itshape} inside an environment are local to it. This
paragraph is italic; the one after the environment is not.
\end{aside}

Back to upright text: the group opened by \verb|\begin{aside}| was closed by
\verb|\end{aside}|, and the font change went with it.

\subsection{Star arguments with xparse}

\cmd{NewDocumentCommand} takes an argument specification instead of a count. Below,
\texttt{s} is an optional star, \texttt{m} a mandatory argument and \texttt{o} an
optional bracketed argument with no default:

\begin{verbatim}
\NewDocumentCommand{\concept}{s m o}{%
  \IfBooleanTF{#1}{\emph{#2}}{\textbf{#2}}%
  \IfValueT{#3}{ (#3)}}
\end{verbatim}
\NewDocumentCommand{\concept}{s m o}{%
  \IfBooleanTF{#1}{\emph{#2}}{\textbf{#2}}%
  \IfValueT{#3}{ (#3)}}

So \verb|\concept{catcode}| prints \concept{catcode}, \verb|\concept*{catcode}| prints
\concept*{catcode}, and \verb|\concept{OMML}[Office Math]| prints
\concept{OMML}[Office Math]. \cmd{NewDocumentEnvironment}, including the \texttt{b}
argument that captures the whole body, is supported as well.

\subsection{Primitive definitions}

Plain \TeX{} definitions work too. A \cmd{def} may use \emph{delimited} parameters, where the
text between parameters must appear literally in the input:

\begin{verbatim}
\def\interval[#1..#2]{$[#1,#2]$}
\end{verbatim}
\def\interval[#1..#2]{$[#1,#2]$}

Here \verb|\interval[0..1]| becomes \interval[0..1] and \verb|\interval[a..b]| becomes
\interval[a..b]: the expander matched the two dots and the closing bracket to split the
arguments. Switches made with \cmd{newif} behave like their \TeX{} counterparts:

\begin{verbatim}
\newif\ifshowhints
\showhintstrue  \ifshowhints Hints are on.\else Hints are off.\fi
\showhintsfalse \ifshowhints Hints are on.\else Hints are off.\fi
\end{verbatim}
\newif\ifshowhints

After \verb|\showhintstrue| the test prints
``\showhintstrue\ifshowhints Hints are on.\else Hints are off.\fi'' and after
\verb|\showhintsfalse| it prints
``\showhintsfalse\ifshowhints Hints are on.\else Hints are off.\fi'' instead. Other primitives
such as \cmd{edef}, \cmd{let}, \cmd{expandafter}, \cmd{csname} and \cmd{ifx} are
available for the same kind of programming.

\begin{tip}
All expansion runs under budgets for expansion steps, nesting depth, token count and
wall-clock time. A runaway recursive macro therefore stops with error E003, which names
the macro stack, instead of hanging the conversion. A command that is not defined at all
is reported as warning W013, and its arguments are kept as text.
\end{tip}

\subsection{Counters and cross-references}

Custom counters take part in the same numbering machinery as sections and equations. The
counter below resets with every section, and its printed form is redefined to show
the section number followed by a capital letter:

\begin{verbatim}
\newcounter{checkpoint}[section]
\renewcommand{\thecheckpoint}{\thesection.\Alph{checkpoint}}
\newcommand{\checkpoint}[1]{\par\noindent
  \refstepcounter{checkpoint}\textbf{Checkpoint \thecheckpoint.} #1\par}
\end{verbatim}
\newcounter{checkpoint}[section]
\renewcommand{\thecheckpoint}{\thesection.\Alph{checkpoint}}
\newcommand{\checkpoint}[1]{\par\noindent
  \refstepcounter{checkpoint}\textbf{Checkpoint \thecheckpoint.} #1\par}

\checkpoint{The document compiles without errors.}\label{cp:macros-compiles}
\checkpoint{Every label is referenced at least once.}\label{cp:macros-labels}

Because \cmd{refstepcounter} sets the current label, \verb|\ref{cp:macros-labels}|
yields \ref{cp:macros-labels} and \verb|\ref{cp:macros-compiles}| yields
\ref{cp:macros-compiles}, the same text that \LaTeX{} prints. Note that the format,
including the section prefix, comes from your redefinition of \cmd{thecheckpoint}.

\subsection{Customised lists}

The \texttt{enumitem} keys are read too. A list can choose its own label format and a
later list can continue where it stopped, with \texttt{resume*} also reusing the first
list's options:

\begin{verbatim}
\begin{enumerate}[label=(\alph*)]
  \item Write the macro.
  \item Use it in the text.
\end{enumerate}
\begin{enumerate}[resume*]
  \item Check the output.
\end{enumerate}
\end{verbatim}

The live version, with a paragraph between the two lists:

\begin{enumerate}[label=(\alph*)]
  \item Write the macro in the preamble or before its first use.\label{it:macros-write}
  \item Use it in the text as often as needed.
\end{enumerate}
Then compile the document and open the result in Word.
\begin{enumerate}[resume*]
  \item Check the output: the numbering continues with the same format.\label{it:macros-check}
\end{enumerate}

Item references follow the label format, so \verb|\ref{it:macros-check}| gives
\ref{it:macros-check} and \verb|\ref{it:macros-write}| gives \ref{it:macros-write}.

\subsection{Limits}

Code inside \cmd{ExplSyntaxOn}\,\dots\,\cmd{ExplSyntaxOff} is not interpreted; the block
is skipped with warning W015. A package that \texdocx{} does not know is reported as
W010. Packages that are accepted without a warning do not
necessarily implement every one of their commands, and an unknown command still produces
W013. The commands that \texdocx{} recognises are listed in
\texttt{docs/supported-commands.md}. Section~\ref{sec:features} gives an overview of the
other document features, and Section~\ref{sec:security} explains the sandbox in which all
of this expansion runs.
